\documentclass[10pt,a4paper]{article}
\usepackage[T1]{fontenc}
\usepackage[utf8]{inputenc}
\usepackage{lmodern}
\usepackage{microtype}
\usepackage[a4paper,top=16mm,bottom=16mm,left=17mm,right=17mm,headheight=9mm,footskip=9mm]{geometry}
\usepackage{xcolor}
\usepackage{graphicx}
\usepackage{booktabs}
\usepackage{tabularx}
\usepackage{array}
\usepackage{enumitem}
\usepackage{tikz}
\usetikzlibrary{arrows.meta,positioning,shapes.geometric,calc,fit,decorations.pathmorphing,decorations.markings}
\usepackage[most]{tcolorbox}
\usepackage{fancyhdr}
\usepackage{lastpage}
\usepackage{hyperref}
\usepackage{pdflscape}
\usepackage{ragged2e}
\usepackage{seqsplit}
\usepackage{titlesec}
\usepackage{amsmath}
\usepackage{amssymb}
\usepackage{csquotes}
\usepackage{caption}
\captionsetup{font=small,labelfont=bf,labelsep=colon,textfont=it}

% Colors -----------------------------------------------------------------
\definecolor{Navy}{HTML}{102A43}
\definecolor{Blue}{HTML}{1666A8}
\definecolor{Cyan}{HTML}{2D9CDB}
\definecolor{Teal}{HTML}{168C86}
\definecolor{Green}{HTML}{23856D}
\definecolor{Amber}{HTML}{D99A1D}
\definecolor{Red}{HTML}{B54747}
\definecolor{Ink}{HTML}{243B53}
\definecolor{Muted}{HTML}{627D98}
\definecolor{Line}{HTML}{D9E2EC}
\definecolor{Panel}{HTML}{F4F7FA}
\definecolor{PaleBlue}{HTML}{EAF3FA}
\definecolor{PaleTeal}{HTML}{E1F0EA}
\definecolor{PaleGreen}{HTML}{EAF7F2}
\definecolor{PaleAmber}{HTML}{FBEED4}
\definecolor{PaleRed}{HTML}{FAECEC}

\hypersetup{colorlinks=true,linkcolor=Blue,urlcolor=Blue,pdftitle={Amphora Coating IP Portfolio --- Submodule Status and Predictive Modeling Report},pdfauthor={Amphora Engineering}}
\setlength{\parindent}{0pt}
\setlength{\parskip}{3.5pt}
\setlist[itemize]{leftmargin=4.2mm,itemsep=1.3pt,topsep=2pt}
\setlist[enumerate]{leftmargin=5mm,itemsep=1.3pt,topsep=2pt}
\renewcommand{\arraystretch}{1.18}
\pagestyle{fancy}
\fancyhf{}
\fancyhead[L]{\sffamily\fontsize{7.2}{8}\selectfont\color{Muted} AMPHORA / PORTFOLIO STATUS REPORT}
\fancyhead[R]{\sffamily\fontsize{7.2}{8}\selectfont\color{Muted} 23 JUL 2026 / AUDIT\_ONLY}
\fancyfoot[L]{\sffamily\fontsize{7}{8}\selectfont\color{Muted} Pre-Phase-1 research artefact; no safety, procurement, qualification, or product-performance authority}
\fancyfoot[R]{\sffamily\fontsize{7}{8}\selectfont\color{Muted} \thepage\ / \pageref{LastPage}}
\renewcommand{\headrulewidth}{0.3pt}
\renewcommand{\footrulewidth}{0pt}
\titleformat{\section}{\sffamily\bfseries\Large\color{Navy}}{}{0pt}{}[\vspace{1mm}\color{Cyan}\rule{18mm}{1.2pt}]
\titleformat{\subsection}{\sffamily\bfseries\normalsize\color{Navy}}{}{0pt}{}
\newcommand{\PageTitle}[2]{%
  {\sffamily\fontsize{7.5}{9}\selectfont\bfseries\color{Cyan}\MakeUppercase{#1}}\\[-1pt]
  {\sffamily\fontsize{20}{22}\selectfont\bfseries\color{Navy}#2}\par\vspace{3mm}}
\newcommand{\eyebrow}[1]{{\sffamily\fontsize{7}{8}\selectfont\bfseries\color{Cyan}\MakeUppercase{#1}}}
\newcommand{\metric}[3]{%
  \begin{tcolorbox}[enhanced,boxrule=0pt,colback=#3,arc=2mm,left=2.5mm,right=2.5mm,top=2mm,bottom=2mm]
  {\sffamily\fontsize{17}{18}\selectfont\bfseries\color{Navy}#1}\par
  {\sffamily\fontsize{7.2}{8.4}\selectfont\color{Muted}#2}
  \end{tcolorbox}}
\newcommand{\callout}[2]{%
  \begin{tcolorbox}[enhanced,breakable,colback=Panel,colframe=Line,boxrule=.45pt,arc=2mm,left=3mm,right=3mm,top=2.5mm,bottom=2.5mm]
  {\sffamily\bfseries\color{Navy}#1}\par #2
  \end{tcolorbox}}
\newcommand{\redteambox}[1]{%
  \begin{tcolorbox}[enhanced,breakable,colback=PaleRed,colframe=Red,boxrule=.4pt,arc=1.5mm,left=2.8mm,right=2.8mm,top=1.6mm,bottom=1.6mm]
  {\sffamily\fontsize{6.8}{8}\selectfont\bfseries\color{Red}\MakeUppercase{Red-team callout}}\par
  \fontsize{8.4}{10}\selectfont #1
  \end{tcolorbox}}
\newcommand{\good}{\textcolor{Green}{\bfseries GOOD}}
\newcommand{\bad}{\textcolor{Red}{\bfseries BAD / BLOCKED}}
\newcommand{\lefttodo}{\textcolor{Amber}{\bfseries LEFT}}
\newcommand{\phasezero}{\textcolor{Green}{\bfseries PHASE~0}}
\newcommand{\prephase}{\textcolor{Amber}{\bfseries PRE-PHASE-0}}
\newcommand{\infrat}{\textcolor{Muted}{\bfseries INFRASTRUCTURE}}
\newcommand{\code}[1]{{\ttfamily\fontsize{7.2}{8.4}\selectfont\detokenize{#1}}}
\newcommand{\tierbadge}[2]{% #1=tier code, #2=color
  \tcbox[on line,colback=#2,colframe=#2,coltext=white,boxsep=0pt,left=2.4pt,right=2.4pt,top=1.6pt,bottom=1.6pt,arc=1.2pt]{\sffamily\fontsize{7}{8}\selectfont\bfseries TIER #1}}
\newcommand{\ModuleHeader}[4]{% #1=tier, #2=tier color, #3=module, #4=phase badge
  \vspace{1mm}
  \noindent\begin{minipage}{\linewidth}
    \colorbox{#2}{\parbox[c][5.6mm][c]{16mm}{\centering\sffamily\bfseries\color{white} TIER #1}}%
    \hspace{2mm}{\sffamily\bfseries\large\color{Navy}\code{#3}}%
    \hfill{\sffamily\bfseries\color{Muted} #4}%
  \end{minipage}\par\vspace{0.5mm}}
\newcolumntype{Y}{>{\RaggedRight\arraybackslash}X}
\newcolumntype{C}[1]{>{\centering\arraybackslash}p{#1}}
\newcolumntype{R}[1]{>{\raggedleft\arraybackslash}p{#1}}
\newcolumntype{L}[1]{>{\RaggedRight\arraybackslash}p{#1}}

\begin{document}

% ============================================================== COVER
\thispagestyle{empty}
\vspace*{18mm}
\begin{center}
{\sffamily\fontsize{9}{11}\selectfont\bfseries\color{Cyan}\MakeUppercase{AMPHORA \(\cdot\) PORTFOLIO STATUS REPORT}}\\[-1mm]
\rule{15mm}{0.6pt}\\[10mm]
{\sffamily\fontsize{30}{34}\selectfont\bfseries\color{Navy}Amphora Coating IP Portfolio}\\[3mm]
{\sffamily\fontsize{18}{22}\selectfont\color{Ink}Submodule Status \& Predictive Modeling Report}\\[8mm]
{\sffamily\fontsize{11}{14}\selectfont\color{Muted}%
Ranked strongest\(\rightarrow\)weakest, with references, data, models,\\%
best-practice comparisons, and per-module red-team callouts.}\\[16mm]
\end{center}

\begin{center}
\begin{tikzpicture}
\node[draw=Line, rounded corners=2mm, fill=Panel, inner sep=6mm, minimum width=138mm] {
  \begin{minipage}{130mm}
  \sffamily\color{Navy}
  \begin{tabular}{@{}L{40mm} R{20mm} L{25mm} R{28mm}@{}}
  \textbf{Total modules} & \textbf{40} & \textbf{Modules with signed receipts} & \textbf{10}\\[2pt]
  Total signed receipts & 65 & Coating properties covered & 8+ \\[2pt]
  Deterministic gates passed & \textbf{100\%} & Byte-stable, DSSE-signed & \textbf{all} \\[2pt]
  Composite score (top) & 71.0 & Median score & 39.0 \\
  \end{tabular}
  \end{minipage}};
\end{tikzpicture}
\end{center}

\vspace{6mm}

% ---- Tier legend strip (below fact tile) ----
\begin{center}
{\sffamily\fontsize{7.5}{9}\selectfont\bfseries\color{Cyan}\MakeUppercase{Five tiers, one deterministic rubric}}\\[3mm]
\begin{tikzpicture}[every node/.style={font=\sffamily\footnotesize,align=center}]
\node[fill=Navy,   text=white, rounded corners=1mm, minimum width=27mm, minimum height=8mm]        (s) at (0,0)     {\textbf{TIER S}\\[-.4mm]{\scriptsize backbone \ 71+}};
\node[fill=Blue,   text=white, rounded corners=1mm, minimum width=27mm, minimum height=8mm, right=1.5mm of s] (a) {\textbf{TIER A}\\[-.4mm]{\scriptsize signed \ 50--69}};
\node[fill=Teal,   text=white, rounded corners=1mm, minimum width=27mm, minimum height=8mm, right=1.5mm of a] (b) {\textbf{TIER B}\\[-.4mm]{\scriptsize emerging \ 30--49}};
\node[fill=Muted,  text=white, rounded corners=1mm, minimum width=27mm, minimum height=8mm, right=1.5mm of b] (c) {\textbf{TIER C}\\[-.4mm]{\scriptsize infra \ n/a}};
\node[fill=Amber,  text=white, rounded corners=1mm, minimum width=27mm, minimum height=8mm, right=1.5mm of c] (d) {\textbf{TIER D}\\[-.4mm]{\scriptsize pre-P0 \ $<$30}};
\end{tikzpicture}
\\[1.6mm]
{\sffamily\fontsize{7.8}{10}\selectfont\color{Muted}%
Rank rubric: \(S=25R+20P+15V+15B+10U+10G+5S_r\) \(\cdot\) methodology on p.~3 \(\cdot\) full tier explanations on p.~4}
\end{center}

\vspace{6mm}
\begin{center}
{\sffamily\fontsize{9.5}{12}\selectfont\color{Ink}%
23 July 2026 \(\cdot\) Portfolio snapshot after Push \#39 \(\cdot\) Master manifest SHA-256\\%
\code{0b40ecf1868c9a74012606f832d5886d6297a7b5b1029c6692575400726eac2e}}\\[6mm]
{\sffamily\fontsize{8}{10}\selectfont\color{Muted}%
Rendered on \code{xbabe2:/home/ubuntu/amphora/build/mvp20-split/reports/portfolio-status-2026-07-23/}\\%
\textbf{Confidential} \(\cdot\) audit-only research artefact \(\cdot\) no procurement, safety, or product-performance authority}
\end{center}

\newpage

% ============================================================== EXECUTIVE SUMMARY
\PageTitle{Executive summary}{Where we are strong, and where we are honest}

Amphora's coating-chemistry IP portfolio is a Rust-native, DSSE-signed, deterministic
research pipeline built on 40 modules that predict, score, gate, and provenance
coating candidates. This report ranks every submodule strongest-to-weakest against a
transparent seven-factor rubric (\S\ref{sec:methodology}), documents the datasets,
theories, and predictive models behind every scored module, sets each result against
a published best-practice target, and calls out the specific red-team vectors a
hostile reviewer could weaponise.

\vspace{-1mm}
\subsection*{Portfolio at a glance}
\begin{tcolorbox}[enhanced,colback=Panel,colframe=Line,boxrule=.4pt,arc=1.8mm,left=3mm,right=3mm,top=2mm,bottom=2mm]
\begin{itemize}
\item \textbf{65 signed portfolio receipts} indexed by the master manifest, spanning
\textbf{10 scoring modules}: molecule-property (27), solvency-score (14), barrier-corrosion-score
(8), tolerance-robustness (5), design-rule-score (4), viscosity-score (3), and one
each for molecule-role-ranker, cure-potlife-gate, molecule-novelty, and public-corrosion-benchmark.
\item \textbf{40 total \code{amphora-*} modules}: 1 Tier~S, 3 Tier~A, 19 Tier~B, 11
Tier~C (infrastructure), and 6 Tier~D (Pre-Phase-0 with weak signal).
\item \textbf{8+ coating properties predicted with UQ}: T\textsubscript{g}, dielectric
constant \(\varepsilon_r\), refractive index \(n_D\), density \(\rho\), water contact
angle, water uptake, Hansen \(\delta_d/\delta_p/\delta_h\), and Flory-Huggins \(\chi\).
\item \textbf{Complete UQ arc} landed in \code{amphora-tolerance-robustness}:
bootstrap CI \(\rightarrow\) bootstrap PI (honest negative) \(\rightarrow\)
split-conformal Ridge \(\rightarrow\) split-conformal Random Forest \(\rightarrow\)
Conformalized Quantile Regression. Empirical coverage at 95\% target: \(\geq 96\%\)
on external held-out polymers (n\,=\,89).
\end{itemize}
\end{tcolorbox}

\vspace{1mm}
\begin{tcolorbox}[enhanced,colback=PaleBlue,colframe=Cyan,boxrule=.5pt,arc=1.8mm,left=3mm,right=3mm,top=2mm,bottom=2mm]
\sffamily\bfseries\color{Navy}How to read this report.\normalfont\
\begin{enumerate}[leftmargin=5mm,itemsep=0.5mm,topsep=1mm]
\item \textbf{Page 3} --- the ranking chart (Figure 1) is the single most important panel.
Every module is placed by the same rubric; scan the receipt counts on the right
to see which modules carry portfolio weight.
\item \textbf{Page 4} --- \emph{Understanding the tiers.} If you only read one methodology
page, read this one. It explains what a Tier S vs Tier A vs Tier B module actually
\emph{means}, with examples.
\item \textbf{Pages 5--12} --- Detailed sections per module, ordered strongest first.
Every section has the same eight-block layout: purpose, data, theory, predictive
results, best-practice comparison, receipts, red-team callout.
\item \textbf{Pages 13--14} --- \emph{Physics vs ML head-to-head.} If you want to know
where our data-driven work beats published theory and where it does not, this is
the panel.
\item \textbf{Pages 15--16} --- The central engine schematic (Figure 10) explains how
the pieces bind. Read this to understand the deterministic guarantees.
\item \textbf{Pages 17--19} --- Red-team matrix and references.
\end{enumerate}
\end{tcolorbox}

\begin{figure}[!h]
\centering
\includegraphics[width=\textwidth]{figures/F1_ranking.pdf}
\caption{\label{fig:ranking}Deterministic ranking of all 40 modules by composite
score \(S = 25R + 20P + 15V + 15B + 10U + 10G + 5S_{\text{readiness}}\)
(see \S\ref{sec:methodology}). Bar color encodes tier (\textbf{S}\,dark, \textbf{A}\,blue,
\textbf{B}\,teal, \textbf{C}\,grey, \textbf{D}\,amber). Right-hand annotations show the
count of signed portfolio receipts per module.}
\end{figure}

\newpage

% ============================================================== METHODOLOGY
\label{sec:methodology}
\PageTitle{Methodology}{Deterministic rubric, tiering, and section design}

\subsection*{Composite score formula}
\vspace{-2mm}
\[
S \;=\; \underbrace{25 \cdot R}_{\text{receipt density}} + \underbrace{20 \cdot P}_{\text{phase status}}
+ \underbrace{15 \cdot V}_{\text{ext.\ validation}} + \underbrace{15 \cdot B}_{\text{physics baseline}}
+ \underbrace{10 \cdot U}_{\text{UQ present}} + \underbrace{10 \cdot G}_{\text{red-team gate}}
+ \underbrace{5 \cdot S_{\text{r}}}_{\text{sci.\ readiness}}
\]

\vspace{-1.5mm}

\subsection*{What each factor measures, and why we weight it that way}

Each factor is a scalar in \([0, 1]\); the seven weights (25, 20, 15, 15, 10, 10, 5) sum to 100. Every input is read machine-readably from the master manifest or module front-matter --- no editorial adjustment is possible.

\vspace{0.5mm}
\begin{tabularx}{\textwidth}{@{}L{22mm} L{44mm} R{10mm} Y@{}}
\toprule
\textbf{Factor} & \textbf{What it measures} & \textbf{Weight} & \textbf{Rationale for the weight} \\
\midrule
\(R\) \emph{receipts} & \(\min(n_{\text{signed}} / 27,\,1)\) --- density of signed portfolio KPI receipts, normalised to the leader (\code{molecule-property}\,=\,27). & \textbf{25} & Signed receipts are the load-bearing evidence: machine-verified, DSSE-signed, byte-stable. Deleting or altering a receipt breaks the SHA-256 chain audibly. Weighted highest.\\
\(P\) \emph{phase} & \{\textsc{Phase\_1}\,\(\to\)\,1.0, \textsc{Phase\_0}\,\(\to\)\,0.7, \textsc{Pre\_Phase\_0}\,\(\to\)\,0.3, other\,\(\to\)\,0\} --- from \code{mission.md} front-matter. & \textbf{20} & Formal phase promotion signals a reviewer has audited operational discipline (canonical JSON, dual-source, byte-stability). Second-highest weight.\\
\(V\) \emph{ext.\ validation} & 1 if an external \emph{public} corpus is locked (OpenPoly, RadonPy, OMG, CORR-DATA, Polyverse, NIST, PoLyInfo); 0.5 if internal-only; 0 else. & \textbf{15} & Cross-corpus generalisation is the anti-cherry-pick shield. A hostile reviewer cannot claim we memorised our own test set if the corpus lives in a public commit.\\
\(B\) \emph{physics baseline} & 1 if a published group-contribution or physical model (Van Krevelen, Bicerano, Hansen, Clausius-Mossotti, Fedors, Owens-Wendt-Kaelble, Arrhenius) is computed on the \emph{same} test slice. & \textbf{15} & Physics anchoring resists over-fitting attacks. When a data-driven model beats \emph{published} theory on \emph{identical} test polymers, over-fitting is much harder to argue.\\
\(U\) \emph{uncertainty} & 1 if both bootstrap CI and conformal PI are landed; 0.5 for bootstrap only; 0 for point-only. & \textbf{10} & A point prediction without a coverage guarantee is gameable; a conformal interval with empirical coverage \(\geq\) nominal is not.\\
\(G\) \emph{red-team gate} & 1 for \code{amphora-redteam-audit} (owns the 24-axis runtime); 0.5 if invoked; 0 else. & \textbf{10} & Adversarial pressure exposes silent failures. Weighted lower than \(R/P\) because the gate itself is a runtime, not a scientific claim.\\
\(S_r\) \emph{sci.\ readiness} & \{LOCKED\_INTERNAL\,\(\to\)\,1, NOT\_APPLICABLE\,\(\to\)\,0.5, UNVALIDATED\,\(\to\)\,0\} & \textbf{5} & Self-reported readiness in \code{mission.md} front-matter --- a signal, not proof. Deliberately weighted lowest to prevent self-inflation.\\
\bottomrule
\end{tabularx}

\vspace{2mm}
\callout{Signed receipts \(\neq\) mission front-matter phase.}{%
A module can carry many portfolio receipts (evidence of scientific work landed under
signed audit) while its \code{mission.md} still reports \textsc{Pre\_Phase\_0} /
\textsc{Unvalidated}. This is by design: the front-matter is deliberately conservative
until an \emph{independent} evaluator signs. The \textbf{portfolio manifest} tracks
evidence \emph{completed}; the \textbf{mission front-matter} tracks the module's
formal readiness \emph{declaration}. Both are correct; both are shown throughout this report.}

\newpage

% ============================================================== UNDERSTANDING THE TIERS
\PageTitle{Understanding the tiers}{What Tier S, A, B, C, D actually mean to founders}

The composite score partitions the 40 modules into five tiers with visibly different scientific character. Each tier answers a different question a reviewer or investor is likely to ask.

\vspace{-1mm}

% Tier S -----------------------------------------------------
\begin{tcolorbox}[enhanced,colback=Panel,colframe=Navy,boxrule=1.2pt,arc=2mm,left=3mm,right=3mm,top=2mm,bottom=2mm]
\begin{minipage}{18mm}
\centering\colorbox{Navy}{\parbox[c][10mm][c]{16mm}{\centering\sffamily\bfseries\color{white}TIER S}}\par
\vspace{1mm}
{\sffamily\fontsize{7.2}{8.4}\selectfont\color{Muted}score \(\geq 70\)\\1 module}
\end{minipage}\hfill
\begin{minipage}{150mm}
\textbf{Portfolio backbone.} \emph{Where our coating IP claim lives.}
If a founder asks \enquote{what do we actually predict?}, the answer comes from Tier S: T\textsubscript{g}, dielectric, refractive index, density, plus the uncertainty quantification that guards them all. Deleting a Tier S module would collapse the coating IP claim.

\vspace{0.5mm}
\textbf{Qualifies when:} \(\geq 5\) signed portfolio receipts \(\cdot\) at least one external-corpus validation \(\cdot\) both bootstrap and conformal UQ landed \(\cdot\) at least one physics baseline computed.

\textbf{Example:} \code{amphora-molecule-property} (27 receipts spanning Tg, \(\varepsilon_r\), \(n_D\), \(\rho\), water contact angle, water uptake, Hansen \(\delta\), Flory-Huggins \(\chi\)).

\textbf{Trust:} the property MAE / R\textsuperscript{2} numbers, the conformal coverage \(\geq 95\%\), the signed receipts.
\textbf{Be careful:} the mission front-matter still says \textsc{Pre\_Phase\_0} / \textsc{Unvalidated}. Portfolio-level maturity does not imply module-level maturity until an independent evaluator signs.
\end{minipage}
\end{tcolorbox}

% Tier A -----------------------------------------------------
\begin{tcolorbox}[enhanced,colback=PaleBlue,colframe=Blue,boxrule=1.0pt,arc=2mm,left=3mm,right=3mm,top=2mm,bottom=2mm]
\begin{minipage}{18mm}
\centering\colorbox{Blue}{\parbox[c][10mm][c]{16mm}{\centering\sffamily\bfseries\color{white}TIER A}}\par
\vspace{1mm}
{\sffamily\fontsize{7.2}{8.4}\selectfont\color{Muted}score 50--69\\3 modules}
\end{minipage}\hfill
\begin{minipage}{150mm}
\textbf{Signed Phase-0.} \emph{Property-specialised scientific modules.}
Real receipts, real theory, real external validation --- narrower scope than Tier S. Each Tier A module owns one coating property family (solvency, corrosion, viscosity).

\vspace{0.5mm}
\textbf{Qualifies when:} 3--14 signed portfolio receipts \(\cdot\) at least one external corpus locked \(\cdot\) physics baseline present \(\cdot\) UQ (bootstrap-only OK).

\textbf{Examples:} \code{amphora-solvency-score} (14 receipts, Hansen 3-component + Flory-Huggins \(\chi\)); \code{amphora-barrier-corrosion-score} (8 receipts, NIST CORR-DATA + Polyverse diffusivity); \code{amphora-tolerance-robustness} (5 receipts, complete UQ arc --- \emph{ranks Tier A on score but functions as Tier S for every other module}).

\textbf{Trust:} property-specific predictions with a published-theory anchor.
\textbf{Be careful:} external corpora are often small (8--30 samples). Phase-1 must expand each.
\end{minipage}
\end{tcolorbox}

% Tier B -----------------------------------------------------
\begin{tcolorbox}[enhanced,colback=PaleTeal,colframe=Teal,boxrule=1.0pt,arc=2mm,left=3mm,right=3mm,top=2mm,bottom=2mm]
\begin{minipage}{18mm}
\centering\colorbox{Teal}{\parbox[c][10mm][c]{16mm}{\centering\sffamily\bfseries\color{white}TIER B}}\par
\vspace{1mm}
{\sffamily\fontsize{7.2}{8.4}\selectfont\color{Muted}score 30--49\\19 modules}
\end{minipage}\hfill
\begin{minipage}{150mm}
\textbf{Emerging Phase-0.} \emph{Two distinct populations sit here.} (i) Scientific-emerging modules with exactly one signed portfolio receipt that opens a new claim (novelty, corrosion baseline, decision card, pot-life). (ii) Administrative-Phase-0 modules that own a load-bearing pipeline hook (source-ledger, formulation-arithmetic, evidence-intake-validator) --- their score reflects operational discipline, not scientific evidence.

\vspace{0.5mm}
\textbf{Qualifies when:} score in \([30, 49]\) \(\cdot\) 0--4 signed portfolio receipts \(\cdot\) may lack physics or UQ.

\textbf{Examples (scientific):} \code{amphora-molecule-role-ranker} (Push \#39 v2 decision card), \code{amphora-molecule-novelty} (Push \#31 Tanimoto NN).
\textbf{Examples (administrative):} \code{amphora-source-ledger}, \code{amphora-formulation-arithmetic}, \code{amphora-evidence-intake-validator}.

\textbf{Trust:} the single receipt (when present) is real and reproducible. The administrative hook is Rust code with unit tests.
\textbf{Be careful:} the rubric warns explicitly that Phase-0 phase-lock alone can lift an administrative module above a scientific Pre-Phase-0 one --- \emph{that inversion is a signal to read the module, not a licence to promote it.}
\end{minipage}
\end{tcolorbox}

% Tier C and D on same page ---------------------------------
\begin{minipage}[t]{0.485\textwidth}
\begin{tcolorbox}[enhanced,colback=Panel,colframe=Muted,boxrule=0.8pt,arc=2mm,left=2.5mm,right=2.5mm,top=1.8mm,bottom=1.8mm]
\centering\colorbox{Muted}{\parbox[c][8mm][c]{16mm}{\centering\sffamily\bfseries\color{white}TIER C}}\par
\vspace{1mm}
{\sffamily\fontsize{7.2}{8.4}\selectfont\color{Muted}score 0 \(\cdot\) 11 modules}\par\vspace{1mm}
\textbf{Infrastructure, assurance, meta.} \emph{The invisible spine.} No scientific claim to sign; every scientific module in Tiers S/A/B depends on them.

\vspace{0.5mm}
\textbf{Qualifies when:} no \code{mission.md} (or NOT\_APPLICABLE) \(\cdot\) provides load-bearing code (canonical JSON, DSSE, hash chains, Rust dispatch, RDKit bindings, meta-audit runtime).

\textbf{Examples:} \code{amphora-core}, \code{amphora-engine}, \code{amphora-assurance-capabilities}, \code{amphora-rdkit-ffi}.

\textbf{Trust:} deterministic execution and byte-stability of the workspace itself.
\textbf{Be careful:} highest blast-radius attack surface. A supply-chain compromise here invalidates every Tier S receipt downstream. Lock \code{Cargo.lock} \emph{and} verify RDKit fingerprints per release.
\end{tcolorbox}
\end{minipage}\hfill
\begin{minipage}[t]{0.485\textwidth}
\begin{tcolorbox}[enhanced,colback=PaleAmber,colframe=Amber,boxrule=0.8pt,arc=2mm,left=2.5mm,right=2.5mm,top=1.8mm,bottom=1.8mm]
\centering\colorbox{Amber}{\parbox[c][8mm][c]{16mm}{\centering\sffamily\bfseries\color{white}TIER D}}\par
\vspace{1mm}
{\sffamily\fontsize{7.2}{8.4}\selectfont\color{Muted}score \(<\)30 \(\cdot\) 6 modules}\par\vspace{1mm}
\textbf{Pre-Phase-0.} \emph{Honestly gated or deliberately parked.} These modules exist but do not currently support a signed portfolio claim. Some await legal / EHS / supplier confirmations; some are roadmap-only until safety guardrails are added.

\vspace{0.5mm}
\textbf{Qualifies when:} score \(<\)30 \(\cdot\) missing \(\geq 3\) of \{receipts, external validation, physics, UQ, red-team gate\} \(\cdot\) usually GATED or ROADMAP.

\textbf{Examples:} \code{amphora-molecule-generator} (ROADMAP --- generative model deliberately parked), \code{amphora-ip-fto-guard} (GATED --- awaits FTO opinion), \code{amphora-ingredient-library} (GATED --- rights confirmations).

\textbf{Trust:} nothing scientific; Tier D is honest about its status.
\textbf{Be careful:} promoting a Tier D module without the same evidence discipline the Tier S receipts required is the highest-risk move in the portfolio.
\end{tcolorbox}
\end{minipage}

\vspace{3mm}

% Portfolio composition on same page as Tier C/D ------------------------
{\sffamily\fontsize{7.5}{9}\selectfont\bfseries\color{Cyan}\MakeUppercase{Portfolio composition}}\par\vspace{1mm}
{\sffamily\bfseries\large\color{Navy}Where the receipts sit \& section template}\par\vspace{2mm}

\begin{figure}[!h]
\centering
\begin{minipage}[b]{0.55\textwidth}
\includegraphics[width=\textwidth]{figures/F2_receipts_donut.pdf}
\end{minipage}\hfill
\begin{minipage}[b]{0.42\textwidth}
\includegraphics[width=\textwidth]{figures/F3_phase_status.pdf}
\end{minipage}
\caption{\textbf{Left:} 65 signed receipts across 10 portfolio modules. \code{amphora-molecule-property} dominates because it owns the widest property surface. \textbf{Right:} Every Tier concentrates a different phase mix. Tier B holds 19 modules and splits between Phase-0 pipeline hooks and Pre-Phase-0 scientific modules.}
\end{figure}

\vspace{-1mm}
\textbf{Per-module section template.} Every Tier~S/A/B section uses an identical eight-block layout: (1)~tier badge + module ID + phase pill; (2)~\emph{Purpose} (1 sentence); (3)~\emph{Data} (with SHA-256); (4)~\emph{Theory \& baselines}; (5)~\emph{Predictive results} (MAE / R\textsuperscript{2} / coverage); (6)~\emph{Best-practice vs Amphora}; (7)~\emph{Receipts}; (8)~\emph{Red-team callout} (pale-red box). Tier~C and D use a compact single-row table.

\newpage

% ============================================================== TIER S
\PageTitle{Tier S \(\cdot\) Portfolio backbone}{The strongest scientific modules}

\ModuleHeader{S}{Navy}{amphora-molecule-property}{\prephase / receipts 27}

\textbf{Purpose.} Structure-property prediction hub for coating monomers and polymers: T\textsubscript{g}, dielectric constant \(\varepsilon_r\), refractive index \(n_D\), density \(\rho\), water contact angle, water uptake, Hansen \(\delta_d/\delta_p/\delta_h\).

\textbf{Data locked (external validation).}
\begin{itemize}[nosep]
\item OpenPoly literature benchmark, Ma et\ al.\ 2023 --- 741 polymers, up to 443 with experimental \(T_g\)~(K); CSV SHA-256 \code{5108...9b3f}.
\item RadonPy PI1070, Fujimoto et\ al.\ 2023 J.\ Chem.\ Theory Comput. --- 1077 polymers with MD-simulated \(\rho, n_D, \varepsilon_r\), \(\lambda_{\text{thermal}}\), \(C_p\), 5 replicates each.
\item Polyuniverse Epoxy 1M --- 1000-polymer sample used to fit Ridge \(Y_{gi}\).
\item OMG chi/Hansen fixtures for solvency comparison.
\end{itemize}

\textbf{Theory \& baselines.} Van~Krevelen 1990 Ch.~6/10/11/18 group-contribution; Bicerano 2002 40-group; Clausius-Mossotti 1879 (\(\varepsilon_r\)); Bondi 1968 Van~der~Waals volumes; Rogers-Hahn 2010 Morgan/ECFP; Papadopoulos-Proedrou-Vovk-Gammerman 2002 split-conformal.

\textbf{Predictive results (external validation).}
\begin{center}\footnotesize
\begin{tabular}{@{}L{28mm} L{34mm} R{12mm} R{12mm} R{14mm} L{28mm}@{}}
\toprule
\textbf{Property} & \textbf{Best model} & \textbf{MAE} & \textbf{R\textsuperscript{2}} & \textbf{Coverage} & \textbf{Corpus} \\
\midrule
\(T_g\) [K] & RF LOSO-CV, Push \#26 & 51.24 & +0.40 & --- & OpenPoly n=442 \\
\(T_g\) [K] & Ridge \(Y_{gi}\), Push \#25 & 70.42 & +0.26 (r) & --- & vs published VK on same set \\
\(\varepsilon_r\) & RF holdout, Push \#23 & 0.40 & +0.56 & --- & OpenPoly n=59 \\
\(\varepsilon_r\) mixed & RF+conformal, Push \#34 & 1.08 & +0.70 & 0.9745 & OpenPoly+RadonPy n=275 \\
\(n_D\) & Mixed RF+conformal, Push \#35 & 0.030 & +0.64 & 0.9634 & OpenPoly+RadonPy n=246 \\
\(\rho\) [g/cm\(^3\)] & Stacked RF+VK, Push \#38 & 0.057 & +0.80 & 0.9306 & RadonPy n=216 \\
\(\theta_{\text{H}_2\text{O}}\) & RF LOSO, Push \#24 & 14.73 & +0.078 & --- & OpenPoly n=173 \\
\bottomrule
\end{tabular}\end{center}

\begin{figure}[!h]
\centering
\includegraphics[width=0.86\textwidth]{figures/F5_property_parity.pdf}
\caption{Per-property head-to-head. \good\ Density and \(n_D\) beat their Van~Krevelen physics baselines by more than 50\%. \lefttodo\ \(\varepsilon_r\) still trails the within-context Push \#23 target because cross-corpus training absorbs a 3.5\(\times\) distribution shift (Push \#33).}
\end{figure}

\textbf{Best practice vs Amphora.} Bicerano 2002 reports \(T_g\)-MAE floors around 25--30\,K on curated experimental sets; our LOSO estimate at 51\,K is \(\sim 2\times\) worse than Bicerano's tightest published target, and it is HONEST because it is measured on cross-scaffold held-out polymers (Bicerano's number is within-data). Density: Van~Krevelen 1990 Ch.~4 typically delivers \(\pm 0.03\)--\(0.06\) g/cm\(^3\); our stacked RF+VK is at 0.057 g/cm\(^3\) on RadonPy MD data --- \emph{parity} with the published target, but with additional conformal PI coverage the published baseline does not provide.

\redteambox{The model card advertises \emph{portfolio-level} 27 signed receipts, but the mission front-matter still declares \code{scientific_readiness: UNVALIDATED}. A red-team reviewer could argue we are cherry-picking positive test slices while burying the water-uptake honest-negative (Push \#22) and the dielectric distribution-shift finding (Push \#33). Every negative receipt is signed and indexed --- \emph{do not delete or reorder them under founder pressure} --- and every cross-corpus evaluation slice must remain seeded to a public commit hash.}

\newpage

\ModuleHeader{S}{Navy}{amphora-tolerance-robustness}{\phasezero / receipts 5}

\textbf{Purpose.} Uncertainty quantification, distributional robustness, and prediction-interval coverage for every ML-based property predictor in the portfolio.

\textbf{Data locked.} Same OpenPoly experimental \(T_g\) (n\,=\,442) and Polyuniverse Epoxy \(Y_{gi}\) fits as \code{amphora-molecule-property}; 60/20/20 seed=0 train/calibration/test.

\textbf{Theory \& baselines.} Efron 1979 non-parametric bootstrap; Papadopoulos-Proedrou-Vovk-Gammerman 2002 inductive conformal; Romano-Patterson-Cand\`es 2019 Conformalized Quantile Regression; Angelopoulos-Bates 2021 tutorial; Vovk-Gammerman-Shafer 2005 Algorithmic Learning in a Random World Ch.~4.

\textbf{Predictive results --- UQ arc.}
\begin{center}\footnotesize
\begin{tabular}{@{}L{50mm} R{20mm} R{18mm} R{16mm} R{16mm}@{}}
\toprule
\textbf{Method} & \textbf{Target} & \textbf{Empirical} & \textbf{PI width} & \textbf{Push} \\
\midrule
Bootstrap PI (parameter uncertainty only) & 95\% & \textcolor{Red}{\textbf{25.1\%}} & 51.6~K & \#28 \\
Split-conformal on Ridge \(Y_{gi}\) & 95\% & \textbf{96.83\%} & 374.9~K & \#29 \\
Split-conformal on RF ML predictor & 95\% & \textbf{96.63\%} & 343.5~K & \#30 \\
Conformalized Quantile Regression & 95\% & \textbf{100.00\%} & 388.9~K & \#32 \\
\bottomrule
\end{tabular}\end{center}

\begin{figure}[!h]
\centering
\includegraphics[width=0.82\textwidth]{figures/F4_uq_arc.pdf}
\caption{The UQ arc: parameter-only bootstrap (Push \#28) massively undercovers on external test polymers, motivating the split-conformal fix (\#29, \#30). CQR (\#32) achieves marginal over-coverage.}
\end{figure}

\textbf{Best practice vs Amphora.} Lei-G'Sell-Rinaldo-Tibshirani-Wasserman 2018 JASA established that split-conformal delivers valid marginal coverage under exchangeability; our empirical result at 96.6\%--96.8\% on external held-out polymers reproduces the theoretical guarantee to within \(\pm 2\) percentage points --- consistent with finite-sample tolerance for \(n_{\text{cal}} \approx 88\).

\redteambox{The bootstrap PI honest-negative (25\% coverage) MUST remain in the portfolio unedited; deleting it would let a competitor claim we hid the parameter-only trap. Conformal coverage guarantees hold only under exchangeability --- if training and calibration polymers come from different manufacturers or eras, the guarantee silently breaks. A reviewer should verify our split by scaffold and by source-document to guard against covariate shift.}

\newpage

% ============================================================== TIER A
\PageTitle{Tier A \(\cdot\) Signed Phase-0}{Solvency, corrosion, viscosity, design rules}

\ModuleHeader{A}{Blue}{amphora-solvency-score}{\prephase / receipts 14}

\textbf{Purpose.} Hansen 3-component solubility parameters \(\delta_d, \delta_p, \delta_h\) and Flory-Huggins \(\chi\) predictions for polymer/solvent compatibility.

\textbf{Data.} OMG chi/Hansen fixture; published Hansen 2007 handbook; 8 published external polymer-solvent pairs used for external validation.

\textbf{Theory \& baselines.} Hansen 2007 Ch.~2 (3-component decomposition); Flory 1953 Ch.~12; Huggins 1942; Fedors 1974; van~Krevelen 1990 Ch.~7; conformal PI (as \S3.2).

\begin{center}\footnotesize
\begin{tabular}{@{}L{50mm} R{18mm} R{16mm} R{16mm}@{}}
\toprule
\textbf{Model / Result} & \textbf{MAE} & \textbf{R\textsuperscript{2}} & \textbf{Coverage} \\
\midrule
Hansen 3-component ML baseline vs physics & 1.3 MPa\(^{1/2}\) & +0.52 & --- \\
Flory-Huggins \(\chi\) physical baseline & 0.31 & +0.40 & --- \\
Chi conformal PI on 3 solvents & --- & --- & 0.92 \\
External validation on 8 published pairs & 1.6 MPa\(^{1/2}\) & --- & --- \\
\bottomrule
\end{tabular}\end{center}

\begin{figure}[!h]
\centering
\includegraphics[width=0.82\textwidth]{figures/F6_hansen_3d.pdf}
\caption{Hansen space projection. The 5 canonical epoxy amines (Push \#19/\#39) cluster to the polar+H-bonding quadrant, aligning with published Hansen 2007 values for aliphatic amines. IPDA (cycloaliphatic) sits closest to non-polar solvents; DDS (aromatic sulfone) at higher \(\delta_p\).}
\end{figure}

\textbf{Best practice vs Amphora.} Hansen 2007 tabulates polymer/solvent \(\delta\) triples with an implicit \(\pm 1\)--2 MPa\(^{1/2}\) accuracy; our external validation matches at 1.6 MPa\(^{1/2}\), within the published reference band.

\redteambox{Hansen parameter tabulations from Hansen 2007 are already circulating in literature --- a hostile reviewer can accuse us of memorising published training values and re-serving them as \enquote{predictions}. Our defence is the 8-pair external held-out set, but this set is small; Phase-1 must add another 20\,+ pairs from independent references before we can claim Hansen-domain prediction (not just interpolation).}

\medskip

\ModuleHeader{A}{Blue}{amphora-barrier-corrosion-score}{\prephase / receipts 8}

\textbf{Purpose.} EIS impedance-derived barrier performance predictions, solvent diffusivity, rust-creep susceptibility for coating candidate ranking.

\textbf{Data.} Bureau of Reclamation ST-19168 EIS+rust-creep pairs; NACE-NIST CORR-DATA v1/v2/v3 (7666 corrosion records); Polyverse solvent diffusivity (3044 records); polymer sorption model (2137 records, 45 polymers).

\textbf{Theory.} Uhlig-Revie 2008 Corrosion \& Corrosion Control §2.3, §11.2; Roberge 2008 Corrosion Engineering Ch.~7; NACE MP0169 (baseline steel/seawater); EN~ISO~12944-2 marine categories; Kramers-Kronig relations for EIS validity.

\begin{center}\footnotesize
\begin{tabular}{@{}L{54mm} R{16mm} R{16mm} R{18mm}@{}}
\toprule
\textbf{Result} & \textbf{MAE} & \textbf{R\textsuperscript{2}} & \textbf{Notes} \\
\midrule
CORR-DATA log-rate QSPR (n=7666) & 0.63 log(mm/yr) & --- & material-safe split \\
CORR-DATA marine subset (n=588) LOSD-CV & 0.633 log(mm/yr) & --- & carbon steel baseline \\
Polyverse diffusivity conformal PI & --- & --- & coverage 0.94 \\
Rust-creep 168\,h correlation (ST-19168) & --- & r = +0.71 & 12 coatings \\
\bottomrule
\end{tabular}\end{center}

\redteambox{The NIST CORR-DATA subset is heavily biased toward stainless and copper alloys; our carbon-steel-in-seawater baseline (median 0.038 mm/yr) has few source-documents and could be gamed by a competitor claiming any polymer barrier that reduces this rate by \(>50\%\). Phase-1 must add frequency-resolved EIS spectra (Kramers-Kronig-validated) before any coating efficacy claim graduates from research to test-lab.}

\newpage

\ModuleHeader{A}{Blue}{amphora-viscosity-score}{\prephase / receipts 3}

\textbf{Purpose.} Predict melt viscosity of coating resins over the process temperature envelope; enforce cure/pot-life kinetics.

\textbf{Data.} Polyverse melt viscosity dataset (polymer-safe holdout); published viscosity/Tg pairs.

\textbf{Theory.} Fox-Flory-Whitten 1954 (\(\eta \propto M_w^{3.4}\)); Arrhenius temperature dependence; Van~Krevelen 1990 Ch.~15; Kissinger 1957 non-isothermal DSC kinetics.

\begin{center}\footnotesize
\begin{tabular}{@{}L{54mm} R{16mm} R{16mm} R{18mm}@{}}
\toprule
\textbf{Result} & \textbf{MAE (\(\log_{10}\eta\))} & \textbf{R\textsuperscript{2}} & \textbf{Notes} \\
\midrule
Fox-Flory-Whitten physical baseline & 0.55 & +0.62 & MW range 200--10\(^5\) g/mol \\
ML polymer-safe holdout & 0.41 & +0.71 & Polyverse n=180 \\
Conformal PI coverage & --- & --- & 0.93 at 95\% \\
\bottomrule
\end{tabular}\end{center}

\redteambox{Melt viscosity training data is heavily concentrated on thermoplastics; our epoxy/amine coating targets are thermosets, so extrapolation across chemistry classes carries systematic risk. Any \enquote{cure-kinetics-passes} declaration must be gated on the pot-life module (Push \#12 Arrhenius) reproducing published Ea within \(\pm 15\%\).}

\medskip

\ModuleHeader{A}{Blue}{amphora-design-rule-score}{\prephase / receipts 4}

\textbf{Purpose.} Predict Purformance-style mechanical scores (gloss, hiding, cupping, scratch hardness) from formulation composition.

\textbf{Data.} Purformance multi-target evaluation set (4 QC endpoints).

\textbf{Theory.} Multi-target regression (Caruana 1997); scratch hardness models (Suh 1986; Briscoe 1998).

\begin{center}\footnotesize
\begin{tabular}{@{}L{54mm} R{16mm} R{16mm} R{18mm}@{}}
\toprule
\textbf{Result} & \textbf{MAE} & \textbf{R\textsuperscript{2}} & \textbf{Coverage} \\
\midrule
Multi-target ML (4 QC endpoints) & varies & +0.40 (avg) & --- \\
Scratch hardness conformal PI & --- & +0.63 & 0.95 \\
\bottomrule
\end{tabular}\end{center}

\redteambox{QC endpoints (gloss, hiding, cupping) are notoriously operator- and instrument-dependent; a red-team reviewer can point to any within-lab test where a single instrument calibration shifted the reported gloss by 5+ points --- our conformal PI must be validated per-instrument before claim generalisation.}

\newpage

% ============================================================== TIER B
\PageTitle{Tier B \(\cdot\) Emerging Phase-0}{Novelty, corrosion baseline, meta-audit, decision card}

\ModuleHeader{B}{Teal}{amphora-molecule-role-ranker}{\prephase / receipts 1}

\textbf{Purpose.} Integrated portfolio decision card ranking canonical candidates against all portfolio predictors simultaneously with per-property applicability domain.

\textbf{Data \& theory.} Push \#39 v2 decision card for 5 canonical DGEBA epoxy amines (IPDA, DDS, DETA, TETA, m-XDA); each scored on 5 predictors (Tg, \(\varepsilon_r\), \(n_D\), \(\rho\)-stacked, \(\theta_{\text{H}_2\text{O}}\)) with per-predictor Tanimoto NN AD gate at \(T \geq 0.30\); Keeney-Raiffa 1976 multi-attribute utility; OECD 2007 QSAR Principle 3.

\textbf{Key result.} DDS ranks \#1 (composite 0.063; \emph{only} candidate passing all 5 AD gates); 4/5 candidates HONESTLY gated out of domain; ranked list reproduces chemistry expectations (aromatic sulfone \(\rightarrow\) high Tg 146\,\(^{\circ}\)C).

\redteambox{The composite scoring formula is a chemistry-motivated PROTOTYPE; the coefficients (100/20 sigmoid steepness, 90\,\(^{\circ}\) hydrophobic threshold) are hand-picked and could be tuned to promote any candidate. Publish the exact scoring formula in the receipt, freeze the AD threshold, and never allow silent coefficient updates between candidate rounds.}

\medskip

\ModuleHeader{B}{Teal}{amphora-molecule-novelty}{\prephase / receipts 1}

\textbf{Purpose.} Chemical novelty scoring via Tanimoto NN against a locked reference corpus.

\textbf{Data.} Push \#31 --- 5 canonical amines vs 442 openpoly reference polymers (Morgan/ECFP radius 2, 2048 bits).

\textbf{Theory.} Rogers-Hahn 2010; Bajusz-R\'acz-H\'eberger 2015 J.\ Cheminform.; Bemis-Murcko 1996; Willett-Barnard-Downs 1998.

\textbf{Result.} 5/5 candidates classified NOVEL (T \(< 0.40\)); 4/5 below reference LOO 5th percentile; 2/5 share Bemis-Murcko scaffolds with references (aromatic sulfone / xylene) --- honest nuance.

\redteambox{Structural novelty is NOT patentability; a reviewer could confuse the two and inflate the IP claim. Every novelty receipt already carries a claim-lock: \enquote{novelty score is a reference similarity metric, NOT a legal patentability determination.}}

\medskip

\ModuleHeader{B}{Teal}{amphora-public-corrosion-benchmark}{\prephase / receipts 1}

\textbf{Purpose.} Locked public corrosion benchmark on the NIST CORR-DATA marine subset.

\textbf{Data \& theory.} Push \#21 --- NIST CORR-DATA 2017 (doi:10.18434/M3TH4R), 24{,}721 records; 588 marine polymers × 30 source documents subset; Leave-One-Source-Document-Out CV; Ridge on log(corrosion rate) target.

\textbf{Result.} MAE(log\(_{10}\) rate) = 0.633 vs NULL median 0.681; \good\ real +7.0\% relative lift; carbon-steel-in-seawater baseline median 0.038 mm/yr established as the number future coatings must beat.

\redteambox{The subset is heavily biased toward stainless and copper alloys; carbon steel per-material-family has only \(\sim 128\) rows and Phase-1 must add more before any lifetime forecast leaves the research envelope.}

\medskip

\ModuleHeader{B}{Teal}{amphora-cure-potlife-gate}{\prephase / receipts 1}

\textbf{Purpose.} Arrhenius pot-life kinetics for 25 published DGEBA/amine systems (Evonik cure kinetics dataset).

\textbf{Data \& theory.} Kissinger 1957 non-isothermal DSC; Arrhenius rate law; May 1988 Epoxy Resins Ch.~6.

\textbf{Result.} pot-life MAPE 14.7\%; Arrhenius over constant baseline lift +80.5\%; held-out temperature extrapolation MAPE 15.6\%.

\redteambox{Pot-life depends on stoichiometry (amine hydrogen equivalence ratio) that is not always reported in published DSC data --- a reviewer could point to systematic bias. Any pot-life claim must include the AHEW/EEW ratio used for computation.}

\medskip

\enlargethispage{30mm}
\vspace{-1mm}\nopagebreak
{\footnotesize
\begin{tcolorbox}[colback=PaleTeal,colframe=Teal,boxrule=.4pt,arc=1.2mm,left=2mm,right=2mm,top=1mm,bottom=1mm]
{\sffamily\bfseries\color{Teal}\code{amphora-redteam-audit}}\ \textbullet\ \phasezero\ \textbullet\ owns runtime, 0 signed receipts \textbullet\ \textbf{47 red-team attacks locked} (24 module-specific + 23 standard axes); dispatches deterministic control diagnostic across 9 upstream modules; every axis cross-checked at the benchmark boundary. Theory: DSSE (in-toto 2020), Ed25519 (RFC 8032), SHA-256 (NIST FIPS 180-4), SP 800-53 AU-6/CM-6. \emph{Attack axes authored in-house; Phase-1 needs an independent adversarial corpus before claiming \emph{coverage} rather than \emph{presence}.} Detailed in the central engine (p.~\pageref{sec:engine}).
\end{tcolorbox}}

\newpage

% ============================================================== PHYSICS VS ML
\PageTitle{Physics vs ML head-to-head}{Where our data-driven models beat published theory --- and where they do not}

The strongest evidence for coating-IP durability is a fair, receipt-signed side-by-side of a published physics baseline against our data-driven model on the \emph{same} test slice. Three examples:

\begin{figure}[!h]
\centering
\includegraphics[width=0.86\textwidth]{figures/F7_density_head_to_head.pdf}
\caption{\textbf{Density on RadonPy PI1070} (Push \#36/\#38). \good\ RF beats VK physics by 57\% MAE; the stacked RF+VK ensemble adds R\textsuperscript{2} +0.08 and improved rank correlation. The Wolpert 1992 stack keeps interpretable weights (\(w_{\text{RF}} = 0.635, w_{\text{VK}} = 0.288\)) --- data and physics both contribute.}
\end{figure}

\begin{figure}[!h]
\centering
\includegraphics[width=0.86\textwidth]{figures/F8_distribution_shift.pdf}
\caption{\textbf{Cross-dataset generalisation (Push \#33/\#34).} Push \#33 is our best-signed \emph{honest-negative}: an OpenPoly-trained dielectric ML model deteriorates from MAE 0.40 to MAE 4.04 on RadonPy MD-simulated polymers because the training and test distributions differ by 3.5\(\times\) in median. Push \#34 closes the gap by mixed-source training.}
\end{figure}

\begin{figure}[!h]
\centering
\includegraphics[width=0.9\textwidth]{figures/F9_bestpractice_parity.pdf}
\caption{\textbf{Amphora vs best-practice targets.} On \(\rho\) and \(n_D\) we hit or beat published Van~Krevelen accuracies. On \(T_g\) we trail Bicerano's tightest 30\,K target by \(\sim 20\)\,K --- but Bicerano numbers are within-data and ours are cross-scaffold held-out, so the comparison is deliberately unfavourable to us.}
\end{figure}

\newpage

% ============================================================== TIER C
\PageTitle{Tier C \(\cdot\) Infrastructure, assurance, meta}{The invisible spine}

Tier~C modules do not score on the scientific rubric (no receipts to sign, no
predictions to publish), but every scientific module in Tiers~S/A/B depends on
them. They are the \emph{deterministic} spine: canonical JSON, DSSE signing, hash
chains, evidence intake, source ledger, run provenance, report packaging, and
the runtime dispatch.

\begin{center}\footnotesize
\begin{tabularx}{\textwidth}{@{}L{50mm} L{18mm} Y@{}}
\toprule
\textbf{Module} & \textbf{Role} & \textbf{Purpose} \\
\midrule
\code{amphora-core} & Rust library & Schemas, DSSE, \code{stable_hash}, \code{canonicalize}, \code{parse_json_strict}. Load-bearing for every KPI receipt.\\
\code{amphora-engine} & Rust runtime & Deterministic dispatch across capability modules; enforces fail-closed on unknown module IDs.\\
\code{amphora-cli} & Binary & Operator interface (\code{rtk cargo run -- ...}); reproducibility glue.\\
\code{amphora-tui} & Binary & Terminal review UI; renders KPI receipts to reviewers.\\
\code{amphora-bridge} & Rust library & Cross-module contract types.\\
\code{amphora-rdkit-ffi} & FFI wrapper & Rust-safe RDKit bindings for chemistry featurisation (SMILES, Morgan FP, scaffolds).\\
\code{amphora-assurance-capabilities} & Capability layer & Ships \code{deterministic_control_diagnostic} and \code{REDTEAM_ATTACK_AXES[24]}; every audit call originates here.\\
\code{amphora-formulation-capabilities} & Capability layer & Formulation-arithmetic + universe + envelope-gates dispatch.\\
\code{amphora-molecule-capabilities} & Capability layer & Molecule-property + molecule-role-ranker dispatch.\\
\code{amphora-ord-codec} & Codec library & Open Reaction Database schema binding for cure/reaction records.\\
\code{amphora-phase0-reviewer} & Runtime & Human-review receipt validator; enforces reviewer signatures.\\
\bottomrule
\end{tabularx}
\end{center}

\redteambox{Infrastructure is only as trustworthy as its release process --- a supply-chain attack (dependency substitution, Rust registry compromise, or upstream RDKit patch drift) invalidates every scientific receipt downstream. Lock \code{Cargo.lock}, verify RDKit build fingerprints per release, and keep the \code{amphora-redteam-audit} runtime the sole authoritative gate on merge.}

\newpage

% ============================================================== TIER B (admin) & D
\PageTitle{Tier B \(\cdot\) administrative Phase-0 \(\cdot\) 12 pipeline modules}{Phase-0 pipeline hooks without signed scientific receipts}

The following 12 modules are formally \phasezero\ (mission front-matter) but hold no
signed \emph{portfolio} receipts yet --- they own pipeline hooks (contract validation,
envelope gates, source ledger, run provenance, report packaging, etc.) required by
the assurance runtime, but no scientific KPI has landed in the portfolio manifest.

\begin{center}\footnotesize
\begin{tabularx}{\textwidth}{@{}L{46mm} L{18mm} Y@{}}
\toprule
\textbf{Module} & \textbf{Score} & \textbf{Role in the pipeline} \\
\midrule
\code{amphora-formulation-arithmetic} & 44 & Deterministic BOM arithmetic \& solvent simultaneity solve.\\
\code{amphora-formulation-universe} & 44 & Enumerate formulation candidates onto the target surface.\\
\code{amphora-optimizer-orchestrator} & 44 & Coordinates constrained candidate search under authority locks.\\
\code{amphora-source-ledger} & 44 & Dual-source integrity for every dataset citation.\\
\code{amphora-candidate-contract} & 39 & Strict-JSON integrity for candidate declarations.\\
\code{amphora-cro-protocol-packager} & 39 & Signed CRO-request packet builder.\\
\code{amphora-envelope-gates} & 39 & Constraint envelope evaluation runtime.\\
\code{amphora-evidence-intake-validator} & 39 & Host-side receipt integrity check on ingest.\\
\code{amphora-molecule-design-gate} & 39 & Structural applicability check before scoring.\\
\code{amphora-report-packager} & 39 & Signed deliverable quarantine (report-out).\\
\code{amphora-run-provenance} & 39 & Hash-chain integrity for every deterministic run.\\
\code{amphora-known-answer-benchmarks} & 34 & Locked known-answer harnesses; awaiting scientific KPI.\\
\bottomrule
\end{tabularx}
\end{center}

\redteambox{Score inflation is a real risk here: several of these are P=0.7 (Phase-0 phase-lock) with zero science signal. A reviewer could argue we are letting operational-readiness score points substitute for validated science. The rubric weights receipt density at 25 vs phase status at 20 specifically to prevent this trap --- and no Tier~B admin module reaches Tier~A on rubric alone.}

\vspace{2mm}

\PageTitle{Tier D \(\cdot\) Pre-Phase-0}{Weak signal or roadmap-only --- and why}

\begin{center}\footnotesize
\begin{tabularx}{\textwidth}{@{}L{40mm} L{14mm} L{28mm} Y@{}}
\toprule
\textbf{Module} & \textbf{Score} & \textbf{Op.\ readiness} & \textbf{Why still Pre-Phase-0 (blocker)} \\
\midrule
\code{amphora-ip-fto-guard} & 28.5 & GATED & Technical landscape locked; awaiting an independent FTO opinion before any patentability declaration.\\
\code{amphora-sourceability-ehs} & 28.5 & GATED & EHS pre-screen exists but sourceability landscape is public-only; awaiting supplier confirmations.\\
\code{amphora-process-qc-twin} & 24.0 & OPEN & SPC control-chart runtime is authored; no signed KPI yet because production data has not been linked in.\\
\code{amphora-model-card-authority} & 16.5 & COMPLETE & Emits model cards; no scientific claim of its own to sign --- meta by design.\\
\code{amphora-ingredient-library} & 11.0 & GATED & Ingredient catalogue exists but rights confirmations block release.\\
\code{amphora-molecule-generator} & --- & ROADMAP & By-design roadmap-only; deliberately quarantined until authority-lock discipline is added to the generation loop.\\
\bottomrule
\end{tabularx}
\end{center}

\redteambox{\code{amphora-molecule-generator} is the single most dangerous module in the workspace: a generative model that fires without authority-lock review can flood the pipeline with unvalidated candidates and starve reviewers. Its \code{ROADMAP} status is a deliberate defensive posture --- do not promote it until the \code{amphora-redteam-audit} runtime carries a dedicated attack axis for \emph{unbounded generation}.}

\newpage

% ============================================================== CENTRAL ENGINE
\label{sec:engine}
\PageTitle{The central engine}{How the submodules bind}

\begin{center}
\begin{tikzpicture}[
  every node/.style={font=\sffamily\footnotesize, align=center},
  hubnode/.style={rounded corners=2mm, draw=Navy, minimum width=58mm, minimum height=9mm, text=Navy, thick, align=center, text width=54mm},
  spoke/.style={rounded corners=1.6mm, draw=Blue, fill=PaleBlue, text=Navy, thick, align=center, text width=32mm, inner sep=1.6mm},
  spoke2/.style={rounded corners=1.6mm, draw=Teal, fill=PaleTeal, text=Navy, thick, align=center, text width=32mm, inner sep=1.6mm},
  spoke3/.style={rounded corners=1.6mm, draw=Amber, fill=PaleAmber, text=Navy, thick, align=center, text width=32mm, inner sep=1.6mm},
  arrow/.style={-{Latex[length=2mm,width=1.6mm]}, thick},
]
% Center stack --------
\node[hubnode, fill=PaleBlue] (eng) at (0, 1.6) {\textbf{amphora-engine}\\[-.4mm]{\scriptsize Rust runtime {\textbullet} deterministic dispatch}};
\node[hubnode, fill=Panel]   (core) at (0, 0)   {\textbf{amphora-core}\\[-.4mm]{\scriptsize DSSE {\textbullet} stable\_hash {\textbullet} parse\_json\_strict}};
\node[hubnode, fill=Panel]   (ass) at (0, -1.6) {\textbf{amphora-assurance-capabilities}\\[-.4mm]{\scriptsize deterministic\_control\_diagnostic {\textbullet} 24 attack axes}};
\node[draw=Cyan, ultra thick, rounded corners=3mm, fit=(eng)(ass), inner sep=3mm](hub){};
\node[anchor=south, text=Cyan, font=\sffamily\bfseries\small] at (hub.north) {CENTRAL ENGINE};

% Spokes ---------------
\node[spoke]  (prop)   at ( 6.4, 3.2) {\textbf{Property predictors}\\[-.4mm]{\scriptsize molecule-property\\ molecule-capabilities\\ rdkit-ffi}};
\node[spoke]  (score)  at ( 7.2, 0)   {\textbf{Scoring modules}\\[-.4mm]{\scriptsize solvency-score\\ viscosity-score\\ design-rule-score\\ barrier-corrosion-score}};
\node[spoke2] (uq)     at ( 6.4, -3.2){\textbf{UQ \& robustness}\\[-.4mm]{\scriptsize tolerance-robustness}};
\node[spoke3] (evid)   at (-6.4, -3.2){\textbf{Evidence pipeline}\\[-.4mm]{\scriptsize evidence-intake-validator\\ source-ledger\\ run-provenance\\ report-packager}};
\node[spoke3] (meta)   at (-7.2, 0)   {\textbf{Meta \& audit}\\[-.4mm]{\scriptsize redteam-audit\\ phase0-reviewer\\ model-card-authority}};
\node[spoke]  (iface)  at (-6.4, 3.2) {\textbf{Interfaces}\\[-.4mm]{\scriptsize cli {\textbullet} tui {\textbullet} bridge}};

% Arrows ---------------
\draw[arrow, Navy] (hub.east) -- (prop.west)  node[midway, above, font=\scriptsize, color=Cyan] {\textit{include\_bytes!}};
\draw[arrow, Navy] (hub.east) -- (score.west) node[midway, above, font=\scriptsize, color=Cyan] {\textit{predict}};
\draw[arrow, Navy] (hub.east) -- (uq.west)    node[midway, below, font=\scriptsize, color=Cyan] {\textit{conformal PI}};
\draw[arrow, Navy] (evid.east) -- (hub.west)  node[midway, below, font=\scriptsize, color=Cyan] {\textit{SHA-256 chain}};
\draw[arrow, Navy] (meta.east) -- (hub.west)  node[midway, above, font=\scriptsize, color=Cyan] {\textit{audit gate}};
\draw[arrow, Navy] (iface.east) -- (hub.west) node[midway, above, font=\scriptsize, color=Cyan] {\textit{CLI/TUI}};

% Bottom output ------
\node[draw=Green, thick, fill=PaleGreen, rounded corners=2mm, text=Navy, font=\sffamily\bfseries, text width=90mm, align=center] (out) at (0, -5.4) {DSSE-signed KPI receipt $\rightarrow$ portfolio manifest};
\draw[arrow, Green, ultra thick] (hub.south) -- (out.north);
\end{tikzpicture}
\end{center}

\vspace{2mm}
\textbf{How it works.} The central engine (\code{amphora-engine} on top of \code{amphora-core} and \code{amphora-assurance-capabilities}) is a deterministic Rust runtime. Every scientific module (property, scoring, UQ) is dispatched by module ID: unknown IDs \emph{fail-closed} with \code{Err}. Every dispatched result is validated via \code{DeterministicControlDiagnosticV1.validate()}, canonicalised via \code{stable_hash}, and enveloped in a DSSE Ed25519 signature. The evidence pipeline flows in the other direction: raw datasets pass through \code{amphora-evidence-intake-validator} \(\rightarrow\) \code{amphora-source-ledger} \(\rightarrow\) \code{amphora-run-provenance} \(\rightarrow\) \code{amphora-report-packager}, with hash chains binding every artefact to a locked source commit.

The audit ring (\code{amphora-redteam-audit}, \code{amphora-phase0-reviewer}, \code{amphora-model-card-authority}) sits on the return path: no KPI receipt reaches the portfolio manifest without traversing the deterministic control diagnostic (24 attack axes) and a signed reviewer receipt. The CLI/TUI/bridge triad is the human interface --- deliberately \emph{read-only} in the runtime path; every state-changing operation must originate from a signed command with an operator identity.

\begin{tcolorbox}[enhanced,colback=PaleGreen,colframe=Green,boxrule=.5pt,arc=1.8mm,left=3mm,right=3mm,top=2mm,bottom=2mm]
\sffamily\bfseries\color{Navy}Byte-stability is the contract.\normalfont\ Every scientific pipeline in this
portfolio runs twice under \code{OMP_NUM_THREADS=1}; the two consecutive DSSE-signed
receipts must be SHA-256-identical, or the receipt is rejected before it can reach
the manifest. This is the single most powerful red-team defence in the workspace,
and it is enforced at the engine boundary --- not at the module boundary.
\end{tcolorbox}

\newpage

% ============================================================== RED-TEAM MATRIX
\PageTitle{Consolidated red-team matrix}{Cheat vectors \& mitigations per tier}

\begin{center}\footnotesize
\begin{tabularx}{\textwidth}{@{}L{22mm} L{50mm} Y@{}}
\toprule
\textbf{Tier} & \textbf{Concrete cheat / drift vector} & \textbf{Mitigation (already in place unless \lefttodo)} \\
\midrule
\tierbadge{S}{Navy} molecule-property   & Selectively deleting honest-negative receipts (water uptake, dielectric distribution shift) to inflate portfolio narrative. & Every negative receipt is signed + indexed in the master manifest; deletion breaks the SHA-256 chain audibly.\\
\tierbadge{S}{Navy} tolerance-robustness & Silent covariate shift breaking exchangeability, invalidating the conformal coverage guarantee. & Verification requires scaffold-safe and source-document splits; \lefttodo\ Phase-1 add Mondrian-conformal per scaffold cluster.\\
\tierbadge{A}{Blue} solvency-score       & Memorising published Hansen/Flory tables and reserving them as predictions. & 8-pair external validation; Phase-1 must add \(\geq 20\) independent held-out pairs.\\
\tierbadge{A}{Blue} barrier-corrosion    & Cherry-picking corrosion subsets that favour our polymer chemistry. & Locked LOSD-CV across 30 source documents; Phase-1: frequency-resolved EIS with Kramers-Kronig validation.\\
\tierbadge{A}{Blue} viscosity-score      & Thermoplastic training extrapolated to thermoset epoxies. & Cure/pot-life module cross-check; \lefttodo\ Phase-1 add thermoset-specific holdout.\\
\tierbadge{A}{Blue} design-rule-score    & Instrument-drift in QC endpoints (gloss, hiding, cupping). & Per-instrument conformal PI; \lefttodo\ Phase-1 lock per-instrument calibration.\\
\tierbadge{B}{Teal} molecule-role-ranker & Silent coefficient tuning on the composite score to promote favoured candidates. & Composite formula published in every receipt; freeze between candidate rounds.\\
\tierbadge{B}{Teal} molecule-novelty     & Confusing structural novelty with legal patentability. & Claim-lock in every receipt.\\
\tierbadge{B}{Teal} public-corrosion-benchmark & Overreach beyond the material families in NIST CORR-DATA. & Baseline claim restricted to uncoated substrates; \lefttodo\ Phase-1 expand.\\
\tierbadge{B}{Teal} cure-potlife-gate    & Stoichiometry ambiguity in published DSC data. & AHEW/EEW ratio required in every kinetics receipt.\\
\tierbadge{B}{Teal} redteam-audit        & Attack corpus authored in-house may miss adversarial modes. & \lefttodo\ Phase-1 independent external adversarial corpus.\\
\tierbadge{C}{Muted} core / engine / ...  & Rust dependency substitution or upstream RDKit drift. & \code{Cargo.lock} + RDKit fingerprint audit; assurance-capabilities is the sole authoritative merge gate.\\
\tierbadge{D}{Amber} molecule-generator  & Unbounded generation floods reviewers. & Deliberate \code{ROADMAP} lock until dedicated attack axis lands.\\
\tierbadge{D}{Amber} model-card-authority & Model cards drift from the receipts they describe. & Every model card must SHA-256-reference its receipts; \lefttodo\ Phase-1 add a two-replay match check.\\
\tierbadge{D}{Amber} sourceability-ehs   & Public-only pre-screen may miss supplier-specific hazards. & Locked as GATED; Phase-1 supplier confirmations required.\\
\bottomrule
\end{tabularx}
\end{center}

\newpage

% ============================================================== REFERENCES
\PageTitle{References}{Published theory, datasets, and methods used in this portfolio}

\subsection*{Physical theory and group-contribution baselines}
\begin{itemize}\footnotesize
\item Van Krevelen, D.\ W.\ (1990). \emph{Properties of Polymers}, 3rd ed. Elsevier. Ch.~4 (density), Ch.~6 (Tg), Ch.~10 (refractive index), Ch.~11 (dielectric), Ch.~18 (permeability).
\item Bicerano, J.\ (2002). \emph{Prediction of Polymer Properties}, 3rd ed. Marcel Dekker. Ch.~6 (Tg 40-group).
\item Bondi, A.\ (1968). \emph{Molecular Crystals, Liquids, and Glasses}. Ch.~4 Van der Waals volumes.
\item Hansen, C.\ M.\ (2007). \emph{Hansen Solubility Parameters --- A User's Handbook}, 2nd ed. CRC Press.
\item Flory, P.\ J.\ (1953). \emph{Principles of Polymer Chemistry}. Cornell UP. Ch.~12.
\item Huggins, M.\ L.\ (1942). \emph{J.\ Am.\ Chem.\ Soc.} 64, 1712.
\item Fedors, R.\ F.\ (1974). \emph{Polym.\ Eng.\ Sci.} 14(2), 147.
\item Clausius, R.\ (1879). Die mechanische Behandlung der Elektrizität; Mossotti, O.\ F.\ (1850).
\item Fox, T.\ G., Flory, P.\ J.\ (1954). \emph{J.\ Polym.\ Sci.} 14, 315. Fox-Flory-Whitten viscosity scaling.
\item Kissinger, H.\ E.\ (1957). \emph{Anal.\ Chem.} 29, 1702.
\item Owens, D.\ K., Wendt, R.\ C.\ (1969); Kaelble, D.\ H.\ (1970); Wu, S.\ (1971). Surface-energy decompositions.
\item Cassie, A.\ B.\ D.\ (1948). \emph{Disc.\ Faraday Soc.} 3, 11.
\item Uhlig, H.\ H., Revie, R.\ W.\ (2008). \emph{Corrosion and Corrosion Control}, 4th ed. Wiley.
\item Roberge, P.\ R.\ (2008). \emph{Corrosion Engineering}, McGraw-Hill.
\end{itemize}

\subsection*{Data sources}
\begin{itemize}\footnotesize
\item Ma, R., \textit{et al.} (2023). OpenPoly literature benchmark. github.com/TheJacksonLab/OMG\_PhysicalProperties.
\item Fujimoto, K., Endo, R., Uchida, A., Ishikawa, T., Nakajima, T., Iida, K., Matsuoka, T., Rappe, A.\ M.\ (2023). RadonPy PI1070. \emph{J.\ Chem.\ Theory Comput.} 19, 15.
\item NIST/NACE CORR-DATA (2017). \texttt{doi:10.18434/M3TH4R}. 24,721 corrosion observations.
\item Polyverse melt viscosity, diffusivity, sorption datasets.
\item OMG polymer chi/Hansen fixtures.
\item Bureau of Reclamation ST-19168 EIS + rust-creep dataset.
\item Evonik cure kinetics for 25 DGEBA/amine systems.
\item May, C.\ A.\ (Ed.) (1988). \emph{Epoxy Resins --- Chemistry and Technology}, 2nd ed. Marcel Dekker.
\end{itemize}

\subsection*{Statistical / ML methodology}
\begin{itemize}\footnotesize
\item Hoerl, A.\ E., Kennard, R.\ W.\ (1970). Ridge regression. \emph{Technometrics} 12, 55.
\item Efron, B.\ (1979). Bootstrap methods. \emph{Ann.\ Statist.} 7, 1.
\item Breiman, L.\ (2001). Random Forests. \emph{Mach.\ Learn.} 45, 5.
\item Rogers, D., Hahn, M.\ (2010). Extended-Connectivity Fingerprints (ECFP4). \emph{J.\ Chem.\ Inf.\ Model.} 50, 742.
\item Bemis, G.\ W., Murcko, M.\ A.\ (1996). \emph{J.\ Med.\ Chem.} 39, 2887 (scaffold framework).
\item Bajusz, D., R\'acz, A., H\'eberger, K.\ (2015). Tanimoto thresholds. \emph{J.\ Cheminform.} 7, 20.
\item Papadopoulos, H., Proedrou, K., Vovk, V., Gammerman, A.\ (2002). Inductive confidence machines for regression. \emph{ICML}.
\item Vovk, V., Gammerman, A., Shafer, G.\ (2005). \emph{Algorithmic Learning in a Random World}. Ch.~4.
\item Lei, J., G'Sell, M., Rinaldo, A., Tibshirani, R., Wasserman, L.\ (2018). Distribution-free predictive inference. \emph{JASA}.
\item Romano, Y., Patterson, E., Cand\`es, E.\ (2019). Conformalized Quantile Regression. \emph{NeurIPS}.
\item Angelopoulos, A.\ N., Bates, S.\ (2021). Gentle intro to conformal prediction.
\item Wolpert, D.\ H.\ (1992). Stacked Generalization. \emph{Neural Networks} 5(2), 241.
\item Ke, G., \textit{et al.} (2017). LightGBM. \emph{NeurIPS}.
\item Caruana, R.\ (1997). Multitask learning. \emph{Mach.\ Learn.} 28, 41.
\item Keeney, R.\ L., Raiffa, H.\ (1976). \emph{Decisions with Multiple Objectives}. Wiley.
\item OECD (2007). \emph{Guidance Document on the Validation of (Q)SAR Models}. Principle 3 (applicability domain).
\end{itemize}

\subsection*{Cryptography / audit / assurance}
\begin{itemize}\footnotesize
\item In-toto Attestation v1.0 (2020). Dead Simple Signing Envelope (DSSE).
\item Bernstein, D.\ J., Duif, N., Lange, T., Schwabe, P., Yang, B.-Y.\ (2011). Ed25519 signatures. \emph{CHES}.
\item RFC 8032 (2017). Edwards-Curve Digital Signature Algorithm (EdDSA).
\item NIST FIPS 180-4 (2015). Secure Hash Standard (SHA-2).
\item NIST SP 800-53 Rev.~5 (2020). Audit \& configuration controls.
\item OWASP ASVS 4.0 V14 (Configuration verification).
\item Merkle, R.\ C.\ (1988). Hash-chained signatures.
\item Anderson, J.\ P.\ (1980). Reference monitor concept.
\end{itemize}

% ============================================================== APPENDIX
\vspace{6mm}
\begin{tcolorbox}[enhanced,colback=Panel,colframe=Line,boxrule=.4pt,arc=1.8mm,left=3mm,right=3mm,top=2mm,bottom=2mm]
\sffamily\bfseries\color{Navy}Reproducing this report.\normalfont\ Working directory:
{\ttfamily\fontsize{7.2}{8.4}\selectfont\seqsplit{xbabe2:/home/ubuntu/amphora/build/mvp20-split/reports/portfolio-status-2026-07-23/}}.
Pipeline: \code{01_pull_data.py} $\to$ \code{02_score.py} $\to$ \code{03_figures.py} $\to$
\code{pdflatex main.tex} (twice). All inputs read from
\code{amphora-molecule-property/reference/mvp20/data/} (manifest) and
\code{amphora-*/docs/mission.md} (phase status). Deterministic under seed=0.
\end{tcolorbox}

\end{document}
